\section{Research Questions}

Conventional Supervised Fine-Tuning (SFT) is a method of adapting a pre-trained language
model using a smaller dataset that contains labelled examples. For SFT, the model
learns to produce the desired outputs from a given input. The one major issue with
SFT is that it learns from fixed ground-truth outputs. This creates a distribution mismatch between
training and inference \parencite{zhao2601self}. On-policy
distillation (OPD) addresses this mismatch by having the student LLM sample its own trajectories
while a teacher LLM provides dense token-level supervision. However, this also has a problem in that
it requires a larger, teacher LLM. To address both of these issues of SFT and OPD, \textcite{zhao2601self}
introduced On-Policy Self-Distillation (OPSD) in which a single LLM acts as both teacher and student with 
different contexts.

Currently, OPSD work has mainly been demonstrated for LLM reasoning. More recently, 
\textcite{yuan2026vision} introduced Vision-OPD which showed that the idea can be transferred
to multimodal models as well. This is done by using a cropped, evidence specific image
for the teacher and a full-image version for the student. 

The following research questions highlight the intent of this proposal to 
further understand the capabilities of OPSD for visual question answering (VQA):

\begin{enumerate}
    \item To what extent does OPSD improve VQA performance compared with conventional SFT?
    \item How does the effectiveness of OPSD vary across vision-language models with different model sizes and architectures?
\end{enumerate}